\begin{editorial}
Ce fragment présente des permutations et réductions locales
de la coupure ordinaire avec des succédents multiples. Dans le
premier arbre, la règle d'implication à gauche conserve la formule
principale dans sa première prémisse, contrairement au schéma de
\Log{G3c} défini dans ce chapitre. La portée du fragment se limite
aux manipulations affichées.
\end{editorial}






Si $!A \ident (!B \lif !C)$, alors $\pi$ se termine par
\[
\AxiomC{}
\RightLabel{$\pi_1'$}
\Deduce$!B, \Gamma \fCenter \Delta, !C$
\RightLabel{\RightR{\lif}}
\UnaryInf$\Gamma \fCenter \Delta, !B \lif !C$
\LeftSubproofLabel{$\pi_1$}
\AxiomC{}
\RightLabel{$\pi_2'$}
\Deduce$!B \lif !C, \Pi \fCenter \Lambda, !B$
\AxiomC{}
\RightLabel{$\pi_2''$}
\Deduce$!C, \Pi \fCenter \Lambda$
\RightLabel{\LeftR{\lif}}
\BinaryInf$!B \lif !C, \Pi \fCenter \Lambda$
\RightSubproofLabel{$\pi_2$}
\RightLabel{\Cut}
\BinaryInf$\Gamma, \Pi \fCenter \Delta, \Lambda$
\DisplayProof
\]
La !!{proof} suivante, qui se termine par \Cut,
\[
\AxiomC{}
\RightLabel{$\pi_1'$}
\Deduce$!B, \Gamma \fCenter \Delta, !C$
\RightLabel{\RightR{\lif}}
\UnaryInf$\Gamma \fCenter \Delta, !B \lif !C$
\LeftSubproofLabel{$\pi_1$}
\AxiomC{}
\RightLabel{$\pi_2'$}
\Deduce$!B \lif !C, \Pi \fCenter \Lambda, !B$
\RightLabel{\Cut}
\BinaryInf$\Gamma, \Pi \fCenter \Delta, \Lambda, !B$
\DisplayProof
\]
a une !!{height} de coupure inférieure à celle de~$\pi$.
L'hypothèse de récurrence fournit une !!{proof}~$\pi_1''$ de
$\Gamma, \Pi \fCenter \Delta, \Lambda, !B$.
Considérons aussi la !!{proof} qui se termine par \Cut :
\[
\AxiomC{}
\RightLabel{$\pi_1'$}
\Deduce$!B, \Gamma \fCenter \Delta, !C$
\AxiomC{}
\RightLabel{$\pi_2''$}
\Deduce$!C, \Pi \fCenter \Lambda$
\RightLabel{\Cut}
\BinaryInf$!B, \Gamma, \Pi \fCenter \Delta, \Lambda$
\DisplayProof
\]
Son rang et sa !!{height} de coupure sont inférieurs au rang
et à la !!{height} de coupure de~$\pi$.
L'hypothèse de récurrence fournirait donc une
\Log{G3c}-!!{proof}~$\pi_2'''$ de
$!B, \Gamma, \Pi \fCenter \Delta, \Lambda$ ; cette désignation
du calcul est celle de la source et reste soumise à la réserve
initiale. Appliquons alors l'hypothèse de récurrence à la
!!{proof} qui se termine par \Cut :
\[
\AxiomC{}
\RightLabel{$\pi_1''$}
\Deduce$\Gamma, \Pi \fCenter \Delta, \Lambda, !B$
\AxiomC{}
\RightLabel{$\pi_2'''$}
\Deduce$!B, \Gamma, \Pi \fCenter \Delta, \Lambda$
\RightLabel{\Cut}
\BinaryInf$\Gamma, \Gamma, \Pi, \Pi \fCenter \Delta, \Delta, \Lambda, \Lambda$
\DisplayProof
\]
Son rang de coupure est $< \cutr{\pi}$.
Nous obtenons ainsi une !!{proof} de
$\Gamma, \Gamma, \Pi, \Pi \fCenter \Delta, \Delta, \Lambda, \Lambda$.
L'application de la \olref[seq][inv]{prop:G3c-cont-adm} donne enfin
la !!{proof} $\pi'$ de $\Gamma, \Pi \fCenter \Delta, \Lambda$.
